\begin{figure}[htbp]
\centering
\begin{tikzpicture}[x=1cm,y=1cm,every node/.style={font=\small},
 bag/.style={draw,rounded corners=1pt,minimum width=1.95cm,minimum height=.58cm,align=center,inner sep=3pt}]
\node[bag] (a) at (0,0) {bag 1: $\{1,2\}$};
\node[bag] (b) at (2.65,0) {bag 2: $\{3,4\}$};
\node[bag] (c) at (5.3,0) {bag 3: $\{5,6\}$};
\node[bag] (d) at (7.95,0) {bag 4: $\{1,2\}$};
\draw[-{Stealth[length=3pt]}] (a.east)--(b.west);
\draw[-{Stealth[length=3pt]}] (b.east)--(c.west);
\draw[-{Stealth[length=3pt]}] (c.east)--(d.west);
\draw (-.97,.38)--(-.97,.56)--(6.28,.56)--(6.28,.38);
\node[font=\footnotesize] at (2.65,.79) {bags 1--3 must use all six colors};
\draw (1.67,-.38)--(1.67,-.56)--(8.93,-.56)--(8.93,-.38);
\node[font=\footnotesize] at (5.3,-.79) {bags 2--4 must use the same six colors};
\end{tikzpicture}
\caption{Why the color pairs must repeat every three bags. The two windows share bags 2 and 3, so bag 4 must use the same pair as bag 1. Arrows show the order of the bags, not directed graph edges. The numbers name six distinct colors.}
\label{fig:palettes}
\end{figure}
